% Appendix 3: the algebra behind the diagnostics.
%
% Created 2026-08-12 by papers/1_method/09_carve_plan.md. It exists because that algebra had no
% owner: the first-order expansion behind the sandwich and the bias was promised in the body as
% "(see Supplementary Information)" and appeared in NO file of the manuscript, the composition
% identity was sitting inside the derivation appendix, and two quantities the body and the Results
% use (the affine-invariant distance, the finite-sample floor of the anisotropy statistic) were
% defined only in a comment or in a `% src` line pointing at a script.
%
% INPUT ORDER: after 09_appendix_members and before 11_appendix_repairs, so it is Appendix 3.
% Nothing hard-codes an appendix number; every cross-reference goes through \ref.

\begin{appendixbox}
\label{app:diagnostics}
\section*{The diagnostics, derived}

The body states the three checks and the two matrices they produce. This appendix carries what is
behind them: the expansion that turns a score into a parameter displacement, the exact sense in
which the two parts of the distortion compose, and the two scalar summaries with the floor below
which neither can be read.

\subsection*{The sandwich and the bias are one first-order expansion}

Both corrections come from the same step. Expanding the total score around the parameters
$\bm{\theta}_0$ the recordings were generated from,
\begin{equation}
  S(\hat{\bm{\theta}}) \;\approx\; S(\bm{\theta}_0)
    + \nabla_{\bm{\theta}} S(\bm{\theta}_0)\,(\hat{\bm{\theta}}-\bm{\theta}_0),
  \label{eq:score-expansion}
\end{equation}
and using $S(\hat{\bm{\theta}})=0$ at a maximum-likelihood estimate, with the empirical Hessian
replaced by its expectation $-\mathbf{H}$,
\begin{equation}
  \hat{\bm{\theta}}-\bm{\theta}_0 \;\approx\; \mathbf{H}^{-1}S(\bm{\theta}_0).
  \label{eq:displacement}
\end{equation}
Taking the covariance of Eq.~\ref{eq:displacement} over independent recordings gives the sandwich,
\begin{equation}
  \bm{\Sigma} \;=\; \mathbf{H}^{-1}\mathbf{J}\,\mathbf{H}^{-1}
   \;=\; \mathbf{H}^{-1/2}\,\mathbf{C}\,\mathbf{H}^{-1/2},
  \label{eq:sandwich}
\end{equation}
and taking its expectation gives the bias, $b=\mathbf{H}^{-1}\mathbb{E}[S]$, which is nonzero
exactly when the mean score is. The two are therefore the same approximation applied to the two
moments of one displacement, which is why a member can fail them separately: $\mathbb{E}[S]$ can
vanish while $\mathrm{Cov}(S)$ departs from $\mathbf{H}$, and at $\bm{\theta}_{\mathrm{pool}}$ it
does so by construction. When $\mathbf{C}=\mathbf{I}$, Eq.~\ref{eq:sandwich} returns
$\mathbf{H}^{-1}$ and the correction is inert. Where $\mathbf{H}$ is singular, every inverse is
taken on its retained eigenspace (Methods).

Both are first order, and the manuscript does not lean on that. The sandwich is verified directly
against the empirical covariance of the estimates over the simulated recordings
(Figure~\ref{fig:clouds}), and against their multivariate distribution through the coverage of the
nominal region, which is what an expansion cannot certify on its own.
% src: theory/macroir/docs/Likelihood_Information_Distortion/supplement_information_distortion_main.tex
% sections 6 and 7 (DCC and DIB), transcribed 2026-08-12 rather than re-derived. That file is the
% origin of the "(see Supplementary Information)" pointer the body used to carry; the Supplementary
% Information it named was never written, and the pointer resolved to nothing from 2026-08-09 until
% this appendix was created.

\subsection*{How the two parts of the distortion compose}

These compose exactly. The exact identity is
\[ \mathbf{C}=\mathbf{K}\,\mathbf{R}\,\mathbf{K}^\top,\qquad \mathbf{K}=\mathbf{H}^{-1/2}\,\mathbf{J}_s^{1/2}, \]
which holds as algebra for any $\mathbf{H}$ and $\mathbf{J}_s$, since substituting the definitions collapses the chain. The factor $\mathbf{K}$ satisfies $\mathbf{K}\mathbf{K}^\top=\mathbf{C}_s$, so it is a legitimate square-root factor of the sample distortion, but it is not the symmetric principal square root, and the form $\mathbf{C}=\mathbf{C}_s^{1/2}\,\mathbf{R}\,\mathbf{C}_s^{1/2}$ holds only where $\mathbf{H}$ and $\mathbf{J}_s$ commute. What composes with no assumption at all is the information volume,
\[ \log\det\mathbf{C}=\log\det\mathbf{C}_s+\log\det\mathbf{R}, \]
an exact identity because the determinant is multiplicative. The per-parameter diagonal, the trace, the eigenvalues and the Frobenius norm do not compose across the split; only the volume does, and only the volume is decomposed in the maps.
% MOVED 2026-08-12 from 08_appendix_derivation.tex, verbatim, not one word edited. It was the only
% block of that appendix that derives nothing about the likelihood: it is a property of the
% diagnostic, and the figure it licenses (Figure 4--figure supplement 1) reads the product of two
% cells against a total, which is the reading the last sentence bounds.

\subsection*{The two scalars, and the distance they are the components of}

With $\lambda_1,\dots,\lambda_d$ the eigenvalues of $\mathbf{C}$ on the retained subspace, the
magnitude $m=e^{\bar e}$ and the anisotropy $a=e^{s}$ of Eq.~\ref{eq:magnitude-anisotropy} are the
exponentials of the mean and the standard deviation of $\log\lambda_i$. They are not two arbitrary
summaries: they are the two orthogonal components of the affine-invariant Riemannian distance from
the identity,
\begin{equation}
  D_{\mathrm{AI}} \;=\; \Bigl(\textstyle\sum_i \log^2\lambda_i\Bigr)^{1/2}
   \;=\; \sqrt{d\,(\bar e^{\,2}+s^2)},
  \label{eq:dai}
\end{equation}
so that a member sits at $m=a=1$ if and only if $D_{\mathrm{AI}}=0$, and the two coordinates say
whether a departure is a common factor on every direction or a spread between them. The pair is
reported throughout and $D_{\mathrm{AI}}$ is reported nowhere in its place, because the decision the
Results support is exactly the one the pair separates and the single number collapses: whether a
scalar correction can work at all ($a=1$) and what factor it should use ($m$).

Neither reads below its own sampling floor. The anisotropy of an estimated distortion is positive by
construction even when the truth is isotropic, since the eigenvalues of a matrix estimated from a
finite cloud spread on their own. Drawing a thousand estimates in two dimensions from a perfectly
calibrated covariance and forming the same statistic gives a median of $1.038$ and a $95$th
percentile of $1.081$, so a measured spread of that size is not evidence of shape; and that is a
lower bound on the floor, because the reported covariance the statistic is formed against is itself
estimated.
% src: papers/1_method/decisions/recompute/fig2_shape_and_floor.R (20,000 null draws at p=2,
% n=1000; median 1.038, 95th percentile 1.081). The Results quote the 1.081 as "a finite-sample
% floor of at least 1.08" and, before this appendix, its only trace in the manuscript was that
% sentence and a `% src` line: a reader had the number and no way to reconstruct it.

\end{appendixbox}
